\section{Trabajo en equipo y evolución continua}

\subsection{Matriz de roles y responsabilidades}

Un sistema exitoso de aprendizaje por imitación necesita la colaboración de cuatro equipos: alignment, ops, data y legal. La siguiente matriz enumera el output y el checkpoint de cada rol:

\begin{table}[H]
\centering
\begin{tabular}{p{0.22\textwidth} p{0.28\textwidth} p{0.28\textwidth} p{0.1\textwidth}}
\toprule
Rol & Entregable & Checkpoint principal & Frecuencia \\
\midrule
Alignment Engineer & Evidence matrix, policy checkpoint & KL drift < θ, preference score & per model release \\
Ops/SRE & Drift dashboard, latency log & alert resolution time & daily \\
Data Engineer & Demo catalog, metadata & coverage report & weekly \\
Legal/Compliance & Guardrail doc, review minutes & high-risk prompt gating & per sprint \\
\bottomrule
\end{tabular}
\caption{Matriz de responsabilidades del equipo de aprendizaje por imitación}
\end{table}

\begin{warningbox}{No difuminar la responsabilidad}
Si un guardrail determinado no es responsabilidad ni de ops ni de compliance, nadie responderá durante un incident. Hay que usar la matriz para definir con claridad al owner e incorporar el proceso al team handbook.
\end{warningbox}

\subsection{Evolución continua y direcciones de investigación}

En la vía de research, el aprendizaje por imitación puede avanzar hacia sim-to-real, multi-agent collaboration y la inferencia automatizada de reward. La diapositiva 02-11 destaca, dentro de los escenarios multi-agent, el cooperative imitation, el competitive imitation y el dynamic gating.

\begin{importantbox}{Lo esencial de los future experiments}
Cada research experiment debe ir acompañado de una evaluation matrix que confirme, al terminar el experimento, si las métricas vuelven al baseline; si el drift de las métricas es severo, hay que posponer el rollout a producción.
\end{importantbox}

\subsection{Resumen de la sección}

El trabajo en equipo y la research planning forman el doble eje de la evolución continua: la matriz de responsabilidades garantiza que en cada release haya alguien revisando los datos, y las direcciones de investigación, mediante structured experiments, hacen que los pipelines de aprendizaje por imitación sean robustos frente a los futuros escenarios multi-agent y sim-to-real.
